% Source: https://artofproblemsolving.com/community/c3h2228170_cyclic_polygons
% Stable identity: aops-topic-2228170-q01
Let $A_1A_2\ldots A_n$ be a regular polygon with side $s$ and interior angle $\theta$. Let $O$ be the intersection of the perpendicular bisectors of $A_1A_2$ and $A_2A_3$ (they meet, because the adjacent sides $A_1A_2$ and $A_2A_3$ are not parallel). Then $OA_1=OA_2=OA_3$.

\begin{asy}
usepackage("amsmath");
size(5cm);
defaultpen(fontsize(10pt));
int n=7; pair[] V; for(int k=0;k<n;++k) V.push(dir(90+360*k/n));
for(int k=0;k<n;++k) draw(V[k]--V[(k+1)%n],linewidth(0.9));
draw(unitcircle,blue+dashed);
pair O=(0,0);
for(int k=0;k<4;++k) draw(O--V[k],red);
pair m1=(V[0]+V[1])/2, m2=(V[1]+V[2])/2;
draw(O--m1,gray+dashed); draw(O--m2,gray+dashed);
dot(O); label("$O$",O,SE);
label("$A_1$",V[0],N); label("$A_2$",V[1],NW); label("$A_3$",V[2],W); label("$A_4$",V[3],SW);
\end{asy}

Triangles $OA_1A_2$ and $OA_3A_2$ are congruent (SSS: $OA_1=OA_3$, $OA_2$ common, $A_1A_2=A_2A_3=s$), so $OA_2$ bisects the angle at $A_2$: $\angle OA_2A_3=\dfrac{\theta}{2}$. Triangle $OA_2A_3$ is isosceles, so also $\angle OA_3A_2=\dfrac{\theta}{2}$, and therefore
\[\angle OA_3A_4=\theta-\dfrac{\theta}{2}=\dfrac{\theta}{2}=\angle OA_3A_2 .\]
Now triangles $OA_3A_2$ and $OA_3A_4$ are congruent (SAS: $OA_3$ common, $A_3A_2=A_3A_4=s$, equal included angles), so $OA_4=OA_2$. The same argument, applied at $A_4$, $A_5$, \ldots, shows step by step that $OA_5=OA_6=\cdots=OA_n=OA_2$. Hence all vertices are at the same distance from $O$, and they lie on the circle with centre $O$ and radius $OA_1$. $\blacksquare$
